% from ega4-2-fr.tex
\oldpage[IV-2]{168}
à $U_{R_1}(X)$; mais inversement, si $x\in Z$ n'appartient à aucun des
$Z_\lambda$, alors, pour toute générisation $x'$ de $x$, ou bien $x'$
appartient à $U$, ou bien on a $\dim(\mathcal O_{X,x'})\geq2$, ou bien on a
$\dim(\mathcal O_{X,x'})=1$, et comme $x'$ n'appartient à aucun des
$Z_\lambda$, $\overline{\{x'\}}$ est nécessairement une composante
irréductible de $Z$, de codimension 1 dans $X$, distincte des $Z_\lambda$,
donc $\mathcal O_{X,x'}$ est régulier par définition. On en conclut que
$U_{R_1}(X)=X-\bigcup_\lambda Z_\lambda$; comme l'ensemble des $Z_\lambda$
est localement fini dans $Z$, $\bigcup_\lambda Z_\lambda$ est fermé, ce qui
achève de prouver que $U_{R_1}(X)$ est ouvert dans $X$.
\end{proof}

\subsection{\texorpdfstring{Critères pour que $\operatorname{Nor}(X)$ soit ouvert}{Critères pour que Nor(X) soit ouvert}}
\label{IV.2.s6.13-fr}

\begin{env}[6.13.1]
\label{IV.2.s6.13.1-fr}
Étant donné un préschéma localement noethérien $X$, nous noterons
$\operatorname{Nor}(X)$ l'ensemble des $x\in X$ où $X$ est normal; cet
ensemble contient $\operatorname{Reg}(X)$ et est contenu dans l'ensemble
(ouvert) des points $x$ tels que $\mathcal O_{X,x}$ soit intègre (i.e.
l'ensemble des points $x$ où $X$ est réduit, et qui n'appartiennent qu'à
une seule composante irréductible de $X$).
\end{env}

\begin{proposition}[6.13.2]
\label{IV.2.s6.13.2-fr}
Soient $X$ un préschéma localement noethérien réduit, $X'$ son normalisé
(II, 6.3.8). Si le morphisme canonique $f:X'\to X$ est fini,
$\operatorname{Nor}(X)$ est ouvert dans $X$.
\end{proposition}

\begin{proof}
En effet, on a $X'=\operatorname{Spec}(f_*(\mathcal O_{X'}))$ et
$f_*(\mathcal O_{X'})$ est une $\mathcal O_X$-Algèbre cohérente par
hypothèse (II, 6.1.3). Dire que $X$ est normal en un point $x$ signifie que
l'homomorphisme canonique
$\mathcal O_X\to(f_*(\mathcal O_{X'}))_x$ est bijectif (II, 6.3.4); mais
l'ensemble de ces points est ouvert puisque $\mathcal O_X$ et
$f_*(\mathcal O_{X'})$ sont cohérents (0$_{\rm I}$, 5.2.7).
\end{proof}

\begin{corollary}[6.13.3]
\label{IV.2.s6.13.3-fr}
Si $A$ est un anneau intègre noethérien japonais,
$\operatorname{Nor}(\operatorname{Spec}(A))$ est ouvert dans
$\operatorname{Spec}(A)$.
\end{corollary}

\begin{proposition}[6.13.4]
\label{IV.2.s6.13.4-fr}
Soit $X$ un préschéma localement noethérien. Pour que
$\operatorname{Nor}(X)$ soit ouvert, il faut et il suffit que tout point
maximal $x$ de $X$ tel que $x\in\operatorname{Nor}(X)$ (ce qui signifie que
$X$ est réduit au point $x$) soit intérieur à $\operatorname{Nor}(X)$.
\end{proposition}

\begin{proof}
Il n'y a à prouver que la suffisance de la condition. En vertu de la
remarque faite dans (6.13.1), on peut se borner (en remplaçant $X$ par un
ouvert de $X$ dans lequel $X$ est réduit) au cas où $X$ est réduit, i.e.
supposer que les points maximaux de $X$ appartiennent à $U_{S_1}(X)$ et à
$U_{R_1}(X)$; comme
$\operatorname{Nor}(X)=U_{S_1}(X)\cap U_{R_1}(X)$ en vertu du critère de
Serre (5.8.6), on déduit donc de (6.11.7) et (6.12.10) que ces deux ensembles
sont ouverts; par suite $\operatorname{Nor}(X)$ est ouvert dans $X$.
\end{proof}

\begin{corollary}[6.13.5]
\label{IV.2.s6.13.5-fr}
Si $X$ est tel que $\operatorname{Reg}(X)$ soit ouvert dans $X$, alors
$\operatorname{Nor}(X)$ est ouvert dans $X$.
\end{corollary}

\begin{proof}
En effet, tout point maximal $x$ où $X$ est réduit (donc
$\mathcal O_{X,x}$ un corps) appartient à $\operatorname{Reg}(X)$, donc est
intérieur à $\operatorname{Reg}(X)$ par hypothèse, et a fortiori est
intérieur à $\operatorname{Nor}(X)$.
\end{proof}

\begin{proposition}[6.13.6]
\label{IV.2.s6.13.6-fr}
\textnormal{(Nagata).} --- Soient $A$ un anneau intègre noethérien, $K$ son
corps des fractions, $K'$ une extension finie de $K$, $A'$ la fermeture
intégrale de $A$ dans $K'$. Pour que $A'$ soit une $A$-algèbre finie, il
faut et il suffit que les deux conditions suivantes soient satisfaites:
\begin{enumerate}
\item[(i)] Il existe $f\neq0$ dans $A$ tel que la fermeture intégrale $A'_f$
de l'anneau de fractions $A_f$ dans $K'$ soit une $A_f$-algèbre finie;
\item[(ii)] Pour tout idéal premier $\mathfrak p\in\operatorname{Spec}(A)$,
la fermeture intégrale $A'_\mathfrak p$ de l'anneau local $A_\mathfrak p$
dans $K'$ est une $A_\mathfrak p$-algèbre finie.
\end{enumerate}
\end{proposition}

\begin{proof}
Les conditions sont nécessaires, car pour toute partie multiplicative $S$
de $A$, $S^{-1}A'$ est la fermeture intégrale de $S^{-1}A$ dans $K'$ et est
par hypothèse un $S^{-1}A$-module de type fini. Pour voir que les conditions
sont suffisantes, considérons la famille filtrante croissante $(B_\lambda)$
des sous-$A$-algèbres de $K'$ qui sont des $A$-algèbres finies et ont $K'$
pour corps des fractions. Posons $Y=\operatorname{Spec}(A)$,
$X_\lambda=\operatorname{Spec}(B_\lambda)$, désignons par $u_\lambda$ le
morphisme $X_\lambda\to Y$ et soient
$S_\lambda=X_\lambda-\operatorname{Nor}(X_\lambda)$,
$T_\lambda=u_\lambda(S_\lambda)$. On peut se borner aux $B_\lambda$ qui
contiennent un ensemble fini dont l'image dans $A'_f$ est un système de
générateurs du $A_f$-module $A'_f$; cela entraîne, en vertu de l'hypothèse
(i), que $(B_\lambda)_f=A'_f$ pour tout $\lambda$, ou encore que
$u_\lambda^{-1}(D(f))$ est contenu dans $\operatorname{Nor}(X_\lambda)$. En
vertu de (6.13.2), $S_\lambda$ est donc fermé dans $X_\lambda$, et comme
$u_\lambda$ est un morphisme fini, $T_\lambda$ est fermé dans $Y$. Mais
pour tout $\mathfrak p\in\operatorname{Spec}(A)$, il existe en vertu de
(ii) un $\lambda$ tel que $(B_\lambda)_\mathfrak p=A'_\mathfrak p$, et par
suite tous les points
